\section{Introduction}
\label{sec:intro}

Large language models write fluently, yet their prose is often flat and generic---polished
sentences that read as if written for no one in particular. A common diagnosis is that the model
fails to model its reader: it produces text that is locally correct but pitched at a default,
average audience rather than the specific person who will read it. If this diagnosis is right,
there is an appealing, cheap fix. Rather than fine-tune the model, we could simply ask it, in its
chain-of-thought, to \emph{simulate the reader}---and to do so continuously, roughly once per
sentence, so that each next sentence is chosen against a fresh prediction of how the reader is
reacting: confused, bored, skeptical, or curious. This is an inference-time, prompt-only
intervention that anyone could apply. \emph{Does it work?}

\para{Why this matters.} Writing quality is a pervasive, high-value problem for LLMs used as
assistants, tutors, and technical communicators. If verbalizing an incremental,
per-sentence audience simulation genuinely improved writing, it would be an unusually cheap win:
no training, no reward model, no external listener---just a better prompt. The intuition also has
strong roots. The pragmatics literature shows that speakers who model their listener communicate
more successfully~\citep{andreas2016reasoning,fried2023pragmatics,takmaz2023speaking}, and the
Theory of Mind (\tom) literature shows that default LLM social reasoning is
brittle~\citep{ullman2023large,shapira2023clever}, suggesting that an explicit modeling step might
help.

\para{The gap.} Despite this, the exact intervention has never been tested for writing. Every
prior audience-aware method places the audience model \emph{outside} the generator's own
reasoning: a decode-time gradient~\citep{takmaz2023speaking}, a rerank by a separately trained
listener~\citep{andreas2016reasoning,coderesa2025}, an external
reward~\citep{styleinfusion2023}, or a formal belief update~\citep{crsa2025}. Per-sentence
mental-state simulation \emph{inside} natural-language chain-of-thought exists only for
comprehension QA and dialogue~\citep{perceptom2024,thoughttracing2025,tomagent2025}, never for
open-ended writing quality. Meanwhile, \citet{chakrabarty2025aislop} show that generic reasoning
does \emph{not} improve writing quality---so a gain from an \emph{audience-specific}
chain-of-thought would be a genuinely non-trivial result, and any claimed gain must be isolated
from ``more thinking tokens.''

\para{Our approach.} We run the first direct test of the intervention (\secref{sec:method}). We
build 60 audience-targeted writing tasks---explaining a concept to an 8-year-old, a non-technical
executive, a PhD expert, or a skeptic, plus creative stories with an assigned readership---and
generate each under four conditions that differ \emph{only} in how the model reasons about the
reader: \plain (write directly), \cot (generic ``think step by step,'' a token-matched control),
\simtom (describe the reader once, then write), and \persentence (the intervention: plan
sentence by sentence, and after each sentence simulate how this reader reacts, then choose the
next). Every condition is told exactly who the reader is, so any difference is attributable to
the \emph{reasoning strategy}, not to privileged audience information. We judge only clean final
text with two separate, cross-family models---one role-playing the reader for audience fit, one
scoring overall quality by blind pairwise comparison---and add an adaptation-gap control that
measures whether a method \emph{actually uses} the audience rather than just writing well in
general.

\para{What we find.} The specific hypothesis is refuted, but its underlying intuition is
vindicated. Audience modeling does help---but a single, rich audience model written once up front
(\simtom) is the clear winner on every metric, at lower cost. The proposed per-sentence
simulation is the worst or near-worst condition: it loses pairwise on overall quality to
\emph{all three} baselines (win-rates $0.083$--$0.250$, all $p<0.001$ after Holm correction),
scores lowest on engagement and audience fit, shows \emph{no} better audience adaptation (in fact
the smallest adaptation gap, $6.04$ vs.\ \simtom's $7.00$), does not reduce homogenization, and
is the longest ($170$ words vs.\ $147$) and most token-expensive condition. The ordering
replicates with a completely different generator (\llama). The takeaway is blunt: \emph{``think
about the reader'' helps writing; ``re-simulate the reader every sentence'' hurts it.}

\para{Contributions.}
\begin{itemize}[leftmargin=*,itemsep=0pt,topsep=0pt]
    \item We give the first direct, controlled test of per-sentence audience simulation in
    chain-of-thought for open-ended writing, isolating it from a token-matched generic-reasoning
    control and from privileged audience information.
    \item We show the intervention \emph{backfires}: it loses on overall quality to plain
    generation, generic reasoning, and one-shot audience modeling, and---via an adaptation-gap
    control---adapts to the reader \emph{less}, not more, than a one-shot audience model.
    \item We show the intuition behind it is nonetheless correct: one-shot audience modeling
    (\simtom) is the best condition on audience fit, engagement, adaptation, and diversity, at
    lower token cost, identifying \emph{granularity} as the key design axis.
    \item We replicate the full ordering across two generator families and release all $3{,}876$
    cached generations and judgments for reproducibility.
\end{itemize}

\para{Organization.} \secref{sec:related} positions our work in pragmatics, \tom, and
writing-quality research. \secref{sec:method} describes the tasks, conditions, models, and
metrics. \secref{sec:results} reports the main comparison, the adaptation-gap control,
homogenization and cost, and the cross-family replication. \secref{sec:discussion} interprets
\emph{why} per-sentence simulation hurts and discusses limitations, and \secref{sec:conclusion}
concludes.
